\documentclass[11pt,a4paper]{article}

% ======================================================================================
% CONFIGURACIÓN DE PÁGINA Y TIPOGRAFÍA ESTILO MICROSOFT WORD EJECUTIVO
% ======================================================================================
\usepackage[utf8]{inputenc}
\usepackage[spanish,es-nodecimaldot]{babel}
\usepackage[margin=2.5cm, top=2.5cm, bottom=2.5cm, headheight=15pt]{geometry}

% Tipografía moderna Sans-Serif (estilo Arial / Calibri de Microsoft Word)
\usepackage{helvet}
\renewcommand{\familydefault}{\sfdefault}

% Espaciado de párrafos estilo Word (sin sangría, espacio entre párrafos, interlineado 1.15)
\usepackage{parskip}
\setlength{\parskip}{5.5pt}
\setlength{\parindent}{0pt}
\usepackage{setspace}
\setstretch{1.15}

% Matemáticas y símbolos
\usepackage{amsmath,amssymb,amsfonts}
\usepackage{graphicx}

% Tablas ejecutivas estilo Word
\usepackage{booktabs}
\usepackage{array}
\usepackage{tabularx}
\usepackage{multirow}
\usepackage[table]{xcolor}

% Enlaces y URLs largas (permite cortes limpios sin desbordar márgenes)
\usepackage{xurl}
\usepackage{hyperref}

% Encabezados, pies de página y numeración estilo Word (Página X de Y)
\usepackage{fancyhdr}
\usepackage{lastpage}
\usepackage{caption}
\usepackage{titlesec}
\usepackage{enumitem}
\usepackage{microtype}
\usepackage{float}
\usepackage[most]{tcolorbox}

% Diagramas nativos TikZ
\usepackage{tikz}
\usepackage{pgfplots}
\pgfplotsset{compat=1.18}
\usetikzlibrary{shapes.geometric, arrows.meta, positioning, calc, backgrounds, fit}

% ======================================================================================
% PALETA DE COLORES CORPORATIVOS (MICROSOFT WORD EXECUTIVE & HUAWEI BLUE/RED)
% ======================================================================================
\definecolor{wordblue}{RGB}{31, 78, 121}        % Azul ejecutivo Word (#1F4E79)
\definecolor{wordsteelblue}{RGB}{46, 117, 182}  % Azul acero subtítulos Word (#2E75B6)
\definecolor{worddark}{RGB}{38, 38, 38}         % Texto carbón
\definecolor{tableheader}{RGB}{31, 78, 121}     % Encabezado tabla azul marino
\definecolor{tablezebra}{RGB}{242, 246, 252}    % Cebreado de tabla gris-azul tenue
\definecolor{boxbg}{RGB}{244, 248, 253}        % Fondo de callouts
\definecolor{boxborder}{RGB}{180, 205, 235}    % Borde de callouts
\definecolor{huaweired}{RGB}{207, 10, 44}       % Rojo corporativo Huawei
\definecolor{clinicred}{RGB}{220, 38, 38}       % Alerta alto riesgo
\definecolor{clinicamber}{RGB}{217, 119, 6}     % Riesgo moderado
\definecolor{clinicgreen}{RGB}{5, 150, 105}     % Bajo riesgo fisiológico

% Configuración de Hipervínculos
\hypersetup{
    colorlinks=true,
    linkcolor=wordsteelblue,
    citecolor=wordblue,
    urlcolor=wordsteelblue,
    pdftitle={VITA-ES: Reporte Técnico Ejecutivo},
    pdfauthor={Marjorie Cabrera & Alejandro Ramos | Proyecto VITA-ES}
}

% ======================================================================================
% FORMATO DE TÍTULOS ESTILO MICROSOFT WORD (TÍTULO 1, 2, 3)
% ======================================================================================
\titleformat{\section}
    {\color{wordblue}\LARGE\bfseries}
    {\thesection.}{0.8em}{}
    [{\color{wordsteelblue}\vspace{1pt}\hrule height 0.8pt\vspace{4pt}}]

\titleformat{\subsection}
    {\color{wordsteelblue}\Large\bfseries}
    {\thesubsection.}{0.6em}{}

\titleformat{\subsubsection}
    {\color{wordblue}\normalsize\bfseries}
    {\thesubsubsection.}{0.5em}{}

\titlespacing*{\section}{0pt}{14pt}{6pt}
\titlespacing*{\subsection}{0pt}{10pt}{4pt}
\titlespacing*{\subsubsection}{0pt}{8pt}{3pt}

% ======================================================================================
% ENCABEZADOS Y PIES DE PÁGINA (ESTILO DOCUMENTO FORMAL WORD)
% ======================================================================================
\pagestyle{fancy}
\fancyhf{}
\rhead{\footnotesize \textcolor{gray}{VITA-ES $\bullet$ StartupSTEM 2026 / Huawei ICT Competition}}
\lhead{\footnotesize \textcolor{wordblue}{\textbf{Reporte Técnico Ejecutivo: Plataforma VITA-ES}}}
\lfoot{\footnotesize \textcolor{gray}{Topic 1: Innovative AI on Huawei Cloud $\bullet$ SRS El Salvador}}
\rfoot{\footnotesize \textcolor{worddark}{Página \thepage\ de \pageref{LastPage}}}
\renewcommand{\headrulewidth}{0.4pt}
\renewcommand{\footrulewidth}{0.4pt}

% Listas compactas estilo Word
\setlist[enumerate]{topsep=2pt, itemsep=2pt, parsep=0pt, leftmargin=1.5em}
\setlist[itemize]{topsep=2pt, itemsep=2pt, parsep=0pt, leftmargin=1.5em}

% Cajas de texto destacadas (Callouts estilo Word)
\newtcolorbox{wordcallout}[1][]{
    colback=boxbg,
    colframe=wordsteelblue,
    arc=3pt,
    boxrule=0pt,
    leftrule=3.5pt,
    top=6pt, bottom=6pt, left=10pt, right=10pt,
    fonttitle=\bfseries\color{wordblue},
    #1
}

% Formato de epígrafes de figuras y tablas
\captionsetup{font={small,sf}, labelfont={bf,color=wordblue}, skip=6pt}

\begin{document}

% ======================================================================================
% PORTADA EJECUTIVA ESTILO MICROSOFT WORD
% ======================================================================================
\begin{tcolorbox}[colback=white, colframe=wordblue, arc=4pt, boxrule=1pt, left=12pt, right=12pt, top=12pt, bottom=12pt]
    \begin{center}
        {\color{wordblue}\huge \textbf{VITA-ES: Plataforma Inteligente de Triaje Mamario Oportuno mediante Espectroscopía de Bioimpedancia y Fusión Multimodal}} \\[0.3cm]
        {\color{wordsteelblue}\large \textbf{Reporte Técnico: Integración del Dispositivo de Adquisición Portátil con el Motor Analítico de IA en el Borde y Huawei Cloud}} \\[0.2cm]
        {\normalsize \textit{StartupSTEM 2026 $\bullet$ Huawei ICT Competition 2026--2027 (Topic 1: Innovative AI Powered by Huawei Cloud)}}
    \end{center}
    
    \vspace{0.2cm}
    \hrule height 0.6pt
    \vspace{0.25cm}
    
    \begin{minipage}[t]{0.48\textwidth}
        \footnotesize
        \textbf{Equipo Fundador:} Marjorie Cabrera \& Alejandro Ramos \\
        \textbf{Especialidad:} Ingeniería Biomédica (5.º año) \\
        \textbf{Experiencia Regulatoria:} Superintendencia de Regulación Sanitaria (SRS) de El Salvador (Vigilancia y Tecnovigilancia) \\
        \textbf{Hardware Clínico:} Sonda VITA-PRB-0428 (\$300 USD) + Parche Descartable (\$0.90 USD)
    \end{minipage}
    \hfill
    \begin{minipage}[t]{0.48\textwidth}
        \footnotesize
        \textbf{Edge AI:} MindSpore Lite (Ascend 310, $4.2\text{ ms}$, Offline) \\
        \textbf{Cloud AI:} Huawei ModelArts (Ascend 910 Pro, MLOps) \\
        \textbf{Normativa Local:} RTS 11.03.02:21 (Dispositivo Médico Clase II) \\
        \textbf{Fecha y Versión:} Octubre 2026 $\bullet$ Versión 2.5.0-Executive \\
        \textbf{Repositorio Oficial:} \url{https://github.com/Laroguee/vita-es-screening-o-VITA-ES-DEMO.git}
    \end{minipage}
\end{tcolorbox}

\vspace{0.3cm}

% ======================================================================================
% RESUMEN EJECUTIVO (WORD CALLOUT BOX)
% ======================================================================================
\begin{wordcallout}[title=Resumen Ejecutivo y Declaración de Complementariedad]
En El Salvador, solo existen \textbf{18 mamógrafos públicos para atender a 3.2 millones de mujeres}, lo que genera listas de espera de más de 6 meses y provoca que los tumores se diagnostiquen tardíamente con un tamaño promedio de \textbf{4.8~cm}. El 40.1\% de las mujeres rurales se encuentra totalmente marginado del tamizaje. \textbf{VITA-ES} resuelve esta crisis mediante un dispositivo portátil de bioimpedancia eléctrica (\$300 USD) de uso comunitario y un parche de contacto bioeléctrico desechable (\$0.90 USD), operado por personal de enfermería en Unidades Comunitarias de Salud Familiar (UCSF). 

El presente reporte establece formalmente cómo la \textbf{plataforma analítica y de software desarrollada en este repositorio complementa la adquisición física de la señal}. Mientras la sonda de hardware inyecta microcorrientes alternas inocuas y digitaliza la respuesta dieléctrica tisular, la plataforma de software actúa como el \textbf{motor de inteligencia clínica, control de calidad y triaje en 6 etapas}: (1) Validación en tiempo real del contacto piel-electrodo mediante el índice SQI (\textit{Signal Quality Index}); (2) Extracción biofísica automatizada de parámetros de dispersión Cole-Cole; (3) Fusión Multimodal Bayesiana Tardía combinando bioimpedancia (70\%), biomarcadores metabólicos (15\%) y hallazgos físicos (15\%); (4) Model Zoo optimizado con XGBoost y VITA-Net 1D en MindSpore Lite (Ascend 310) que alcanza \textbf{95.24\% de sensibilidad en carcinoma con latencia de 4.2~ms}; (5) Cuantificación de incertidumbre diagnóstica mediante el índice CUI para evitar descartes erróneos en casos atípicos; y (6) Sincronización segura con Huawei Cloud (ModelArts, DCS, RDS con cifrado KMS bajo RTS 11.03.04:23) para tecnovigilancia y re-entrenamiento federado continuo.
\end{wordcallout}

\vspace{0.2cm}
\tableofcontents
\vspace{0.4cm}

% ======================================================================================
% SECCIÓN 1: CONTEXTO NACIONAL Y JUSTIFICACIÓN
% ======================================================================================
\section{El Problema de Salud Pública y Justificación en El Salvador}

\subsection{La Realidad del Cáncer de Mama en El Salvador}
El cáncer de mama es la principal causa de morbilidad y mortalidad oncológica femenina en El Salvador. A pesar de que la supervivencia a cinco años supera el 99\% si la lesión es detectada en Estadio I, el sistema de salud salvadoreño enfrenta un colapso estructural:
\begin{enumerate}
    \item \textbf{Déficit Crítico de Infraestructura:} En la red pública del Ministerio de Salud (MINSAL) y el Instituto Salvadoreño del Seguro Social (ISSS) operan únicamente \textbf{18 mamógrafos funcionales} para una población objetivo de \textbf{3.2 millones de mujeres} (1.6 millones en el grupo de tamizaje prioritario de 30 a 59 años).
    \item \textbf{Tiempos de Espera Inaceptables:} La saturación de los centros de tercer nivel (Hospital Nacional de la Mujer y Hospital Oncológico) genera listas de espera de \textbf{6 meses o más} para una mamografía diagnóstica.
    \item \textbf{Diagnóstico en Fases Avanzadas:} A causa de la demora, el tumor se diagnostica cuando ya mide en promedio \textbf{4.8~cm} de diámetro, reduciendo drásticamente la probabilidad de curación e incrementando el costo de tratamientos paliativos.
    \item \textbf{Exclusión Rural Sistemática:} El \textbf{40.1\% de las mujeres en áreas rurales} carecen por completo de acceso físico o económico al tamizaje mamográfico.
\end{enumerate}

\subsection{Estrategia VITA-ES: Descentralización Comunitaria Permanente}
A diferencia de iniciativas de brigadas itinerantes ocasionales que dejan a las comunidades desatendidas tras su partida, VITA-ES establece un \textbf{modelo de adopción permanente en el primer nivel de atención}:
\begin{itemize}
    \item \textbf{El Dispositivo se Queda en la Clínica:} Se instalan equipos permanentes en las 750 Unidades Comunitarias de Salud Familiar (UCSF/ISSS) y en 1,200 clínicas rurales y jornadas de ONGs.
    \item \textbf{Capacitación Rápida:} Se capacita en solo \textbf{20 minutos} al personal local (enfermería y promotores de salud) para operar el sistema.
    \item \textbf{Resultados Inmediatos:} Triaje completo, indoloro y sin radiación ionizante con dictamen clínico en \textbf{5 minutos}.
\end{itemize}

% --- CUADRO 1: COMPARATIVA DE ALTERNATIVAS (EXTRAÍDA DE LA DIAPOSITIVA 7) ---
\begin{table}[H]
\centering
\small
\caption{Matriz comparativa de tecnologías de detección mamaria (Fuente: StartupSTEM 2026).}
\label{tab:comparativa}
\renewcommand{\arraystretch}{1.25}
\begin{tabularx}{\textwidth}{p{2.8cm} p{3.2cm} p{3.0cm} p{5.0cm}}
\rowcolor{tableheader}
\textcolor{white}{\textbf{Criterio}} & \textcolor{white}{\textbf{Mamografía Digital}} & \textcolor{white}{\textbf{Autoexploración}} & \textcolor{white}{\textbf{Solución VITA-ES}} \\
\toprule
\textbf{Costo por Prueba} & $\sim$\$50.00 USD (Inaccesible) & Gratuita & \textbf{\$5.00 USD comercial / \$0.00 para la paciente} \\
\rowcolor{tablezebra}
\textbf{Ubicación de Acceso} & Hospitales en ciudades & En el hogar & \textbf{En la comunidad rural (Portátil)} \\
\textbf{Operador} & Médico Radiólogo escaso & La propia paciente & \textbf{Enfermería / Personal de salud local} \\
\rowcolor{tablezebra}
\textbf{Efectividad Clínica} & Alta, pero tardía por colapso & Muy baja sensibilidad & \textbf{Alta sensibilidad: detecta cambios no palpables} \\
\textbf{Tipo de Dictamen} & Diagnóstico final por imagen & Subjetivo y tardío & \textbf{Triaje rápido, objetivo, indoloro y no ionizante} \\
\bottomrule
\end{tabularx}
\end{table}

\subsection{Modelo de Negocio y Sostenibilidad Financiera}
El modelo de negocio está diseñado para no trasladar ningún costo a la usuaria:
\begin{itemize}
    \item \textbf{Costo para la Paciente:} \textbf{\$0.00 USD (100\% Gratuito en la UCSF)}.
    \item \textbf{Contratación Centralizada:} El MINSAL y organizaciones no gubernamentales adquieren paquetes masivos de tamizaje preventivo a un precio institucional de \$3.50 USD por prueba.
    \item \textbf{Estructura de Costos del Dispositivo:} Costo de fabricación del hardware: \textbf{\$300.00 USD}. Costo del parche de electrodos biocompatibles descartable: \textbf{\$0.90 USD}.
    \item \textbf{Margen y Retorno de Inversión (ROI):} Margen neto de \textbf{\$2.60 USD por prueba}. El punto de equilibrio se alcanza a las \textbf{116 pruebas por equipo} (atendiendo a 3 pacientes diarias, el ROI se completa en apenas \textbf{1.3 meses}).
\end{itemize}

% ======================================================================================
% SECCIÓN 2: COMPLEMENTARIEDAD HARDWARE - SOFTWARE
% ======================================================================================
\section{Complementariedad entre la Sonda de Adquisición y la Plataforma de Software}

\subsection{La Necesidad del Eslabón Analítico}
La sonda física VITA-PRB-0428 integra un generador analógico de impedancias multifrecuencia (Analog Devices AD5933) conectado a un parche desechable con electrodos concéntricos en acero quirúrgico 316L. Sin embargo, **la adquisición de hardware por sí sola es ciega desde la perspectiva clínica**:
\begin{itemize}
    \item El sensor físico únicamente adquiere números complejos en forma de parte real e imaginaria de corriente y voltaje: $Z(f) = Z'(f) + j Z''(f)$.
    \item Una enfermera comunitaria en una UCSF rural no puede diagnosticar una lesión mamaria a partir de una tabla cruda de ohmios y nanofaradios.
    \item Errores de colocación física (piel deshidratada o cortocircuitos por exceso de gel) generarían diagnósticos falsos positivos o falsos negativos si no existe una capa digital que supervise la integridad de la señal.
\end{itemize}

Por consiguiente, la plataforma de software alojada en este repositorio actúa como el **cerebro complementario que transforma una medición bioeléctrica cruda en una decisión clínica segura y accionable**.

% --- FIGURA 1: ARQUITECTURA DE COMPLEMENTARIEDAD (TIKZ RESIZED) ---
\begin{figure}[H]
\centering
\resizebox{0.95\linewidth}{!}{
\begin{tikzpicture}[
    node distance=0.85cm,
    hw_box/.style={rectangle, draw=wordblue, fill=wordblue!8, rounded corners=3pt, text centered, font=\small, minimum height=1.1cm},
    sw_box/.style={rectangle, draw=wordsteelblue, fill=boxbg, rounded corners=3pt, text centered, font=\small, minimum height=1.1cm},
    cloud_box/.style={rectangle, draw=huaweired, fill=huaweired!8, rounded corners=3pt, text centered, font=\small, minimum height=1.1cm},
    arrow/.style={-{Latex[length=2.5mm]}, thick, worddark}
]
    % FILA 1: ADQUISICIÓN FÍSICA
    \node (p1) [hw_box, text width=4.5cm] {\textbf{1. Adquisición Física} \\ Parche Descartable (\$0.90) \\ 7 Frecuencias (15 kHz--1 MHz)};
    \node (p2) [hw_box, right=0.9cm of p1, text width=4.5cm] {\textbf{2. Sonda VITA-PRB-0428} \\ AD5933 + ESP32-S3 BLE \\ Medición $Z(\omega) = Z' + jZ''$};
    \node (p3) [hw_box, right=0.9cm of p2, text width=4.5cm] {\textbf{3. Trama Digital Cruda} \\ Tramas binarias BLE 5.0 \\ Envío a la Estación Clínica};

    % FILA 2: PLATAFORMA DE SOFTWARE (ESTE REPOSITORIO)
    \node (s1) [sw_box, below=1.0cm of p1, text width=4.5cm] {\textbf{4. Interlock de Calidad (SQI)} \\ Validación contacto y piel seca \\ Filtro de ruido ($SNR > 14\text{ dB}$)};
    \node (s2) [sw_box, below=1.0cm of p2, text width=4.5cm] {\textbf{5. Motor Cole-Cole} \\ Extracción de 9 descriptores \\ Curvas Nyquist calibradas};
    \node (s3) [sw_box, below=1.0cm of p3, text width=4.5cm] {\textbf{6. Fusión Multimodal Tardía} \\ Bioimpedancia 70\% \\ Coimbra 15\% + WDBC 15\%};

    % FILA 3: DECISIÓN CLÍNICA Y NUBE
    \node (d1) [sw_box, below=1.0cm of s1, text width=4.5cm] {\textbf{7. Edge AI (Ascend 310)} \\ MindSpore Lite (4.2 ms) \\ Sensibilidad 95.24\% en Cáncer};
    \node (d2) [sw_box, below=1.0cm of s2, text width=4.5cm] {\textbf{8. Índice CUI (Incertidumbre)} \\ Entropía Shannon Normalizada \\ Bloquea descartes erróneos};
    \node (c1) [cloud_box, below=1.0cm of s3, text width=4.5cm] {\textbf{9. Huawei Cloud Suite} \\ ModelArts MLOps + RDS KMS \\ Tecnovigilancia RTS 11.03.04};

    % Conexiones
    \draw [arrow] (p1) -- (p2);
    \draw [arrow] (p2) -- (p3);
    \draw [arrow] (p3) -- (s1);
    \draw [arrow] (s1) -- (s2);
    \draw [arrow] (s2) -- (s3);
    \draw [arrow] (s3) -- (d1);
    \draw [arrow] (d1) -- (d2);
    \draw [arrow] (d2) -- (c1);

    % Agrupaciones visuales
    \begin{scope}[on background layer]
        \node [draw=wordblue, fill=wordblue!2, fit=(p1)(p2)(p3), inner sep=0.2cm, rounded corners=4pt, label={[wordblue, font=\bfseries]above:CAPA 1: ADQUISICIÓN DE HARDWARE (SONDA PORTÁTIL EN CLÍNICA)}] {};
        \node [draw=wordsteelblue, fill=wordsteelblue!2, fit=(s1)(s2)(s3)(d1)(d2), inner sep=0.2cm, rounded corners=4pt, label={[wordsteelblue, font=\bfseries]left:\rotatebox{90}{\textbf{PLATAFORMA VITA-ES}}}] {};
        \node [draw=huaweired, fill=huaweired!2, fit=(c1), inner sep=0.2cm, rounded corners=4pt, label={[huaweired, font=\bfseries]below:CAPA 3: HUAWEI CLOUD}] {};
    \end{scope}

\end{tikzpicture}
}
\caption{Cadena de valor técnico-clínica de VITA-ES: Articulación entre la adquisición de señales de hardware, la plataforma de inteligencia en el borde y la gobernanza en Huawei Cloud.}
\label{fig:pipeline_completo}
\end{figure}

\subsection{Las 6 Funciones Complementarias de la Plataforma de Software}
\begin{enumerate}
    \item \textbf{Supervisión de Calidad en Vivo (Signal Quality Index - SQI):} Evalúa de forma no invasiva si los electrodos del parche tienen la adherencia y humedad adecuada sobre la piel de la paciente. Si la resistencia de contacto es excesiva ($R_0 > 4,200\ \Omega$) o se produce un puente de gel ($Z' < 50\ \Omega$), el software bloquea la prueba y le muestra a la enfermera instrucciones claras de colocación.
    \item \textbf{Extracción Analítica de Parámetros de Cole-Cole:} Modela matemáticamente el comportamiento interfacial celular $\beta$ y proyecta los descriptores biomédicos invariantes ($I_0, DR, \text{Max IP}, \text{Area}$) sin requerir intervención del operador.
    \item \textbf{Fusión Multimodal Tardía (Enfoque Holístico):} La sonda adquiere el estado biofísico local (70\%), pero la plataforma permite que la enfermera ingrese datos metabólicos (perfil glucémico de Coimbra, 15\%) y antecedentes de masa palpable (Wisconsin WDBC, 15\%), evitando falsos diagnósticos basados en una sola fuente de datos.
    \item \textbf{Inferencia Edge Ultrarrápida sin Conexión (Offline Ready):} Gracias a la optimización de los modelos en MindSpore Lite para chips Huawei Ascend 310, el triaje se realiza de manera local en \textbf{4.2~ms}, garantizando operatividad en cantones rurales sin señal de internet ni telefonía móvil.
    \item \textbf{Cuantificación de Incertidumbre Médica (CUI):} Calcula la Entropía Normalizada de Shannon en las predicciones. Si el tejido arroja características limítrofes (p. ej. atipia celular en fibroadenomas hipercelulares), el sistema suspende el descarte y emite una orden de confirmación ecográfica o biopsia.
    \item \textbf{Gobernanza y Tecnovigilancia en Huawei Cloud:} Sincroniza los expedientes clínicos en bases de datos PostgreSQL RDS cifradas con KMS AES-256 (cumplimiento HIPAA y GDPR), reportando eventos adversos conforme al **RTS 11.03.04:23** y permitiendo re-entrenar los modelos con datos salvadoreños mediante aprendizaje federado en Huawei ModelArts.
\end{enumerate}

% ======================================================================================
% SECCIÓN 3: FUNDAMENTACIÓN BIOFÍSICA
% ======================================================================================
\section{Fundamentación Biofísica y Modelo de Cole-Cole}

\subsection{Comportamiento Dieléctrico Tisular}
La carcinogénesis mamaria altera profundamente la microarquitectura celular, lo que permite diferenciar tejidos mediante corrientes alternas seguras ($I_{\text{RMS}} < 50\ \mu\text{A}$):
\begin{itemize}
    \item \textbf{Tejido Graso Adiposo Normal:} Células rellenas de lípidos no conductores. Presenta una resistencia extracelular masiva ($R_0 > 2,000\ \Omega$).
    \item \textbf{Tejido Glandular Normal y Lesiones Benignas:} Membranas celulares intactas con capacitancia estándar y resistencia intermedia ($R_0 \approx 230 - 290\ \Omega$).
    \item \textbf{Carcinoma Mamario Maligno:} Proliferación descontrolada con alta densidad celular, atipia nuclear, permeabilidad de membrana elevada, pérdida de uniones estrechas e incremento del agua intra y extracelular. Esto se traduce en una drástica reducción de la resistencia a baja frecuencia ($R_0 \approx 394\ \Omega$) y un marcado aumento en la dispersión $\Delta R$.
\end{itemize}

\subsection{Ecuación Clásica de Cole-Cole y Reactancia Pico (Max IP)}
La impedancia celular compleja $Z(\omega)$ responde a la ecuación analítica formulada por Cole \& Cole \cite{cole1941}:

\begin{equation}
    Z(\omega) = R_\infty + \frac{R_0 - R_\infty}{1 + (j \omega \tau)^\alpha}
    \label{eq:cole}
\end{equation}

En la frecuencia de relajación celular $\omega_c = 1/\tau$, la reactancia capacitiva máxima ($-Z''_{\max}$, catalogada clínicamente como \textbf{Max IP}) responde a la deducción exacta:

\begin{equation}
    -Z''_{\max} = \frac{R_0 - R_\infty}{2} \tan\left(\frac{\alpha \pi}{4}\right)
    \label{eq:max_ip}
\end{equation}

A partir de las 7 frecuencias estándar de Jossinet \cite{jossinet1996, uci_breast_tissue, kaggle_breast_tissue} (15.625~kHz a 1.0~MHz), el motor de procesamiento extrae el vector canónico $\mathbf{x} \in \mathbb{R}^9$:

\begin{equation}
    \mathbf{x} = \Big[ I_0,\; \text{PA500},\; \text{HFS},\; \text{DA},\; \text{Area},\; \text{A/DA},\; \text{Max IP},\; \text{DR},\; P \Big]^T
    \label{eq:vector_x}
\end{equation}

% --- FIGURA 2: DIAGRAMA COLE-COLE EN TIKZ (AJUSTADO PARA NO DESBORDAR) ---
\begin{figure}[H]
\centering
\resizebox{0.82\linewidth}{!}{
\begin{tikzpicture}[scale=1.0]
    \draw[->, thick, worddark] (0,0) -- (9.0,0) node[right, font=\small\bfseries] {Resistencia Real $Z'$ ($\Omega$)};
    \draw[->, thick, worddark] (0,0) -- (0,4.5) node[above, font=\small\bfseries] {Reactancia Capacitiva $-Z''$ ($\Omega$)};

    \def\xc{4.6}
    \def\yc{-0.8}
    \def\rad{3.8}

    \draw[very thick, wordblue, domain=35:145, samples=100] 
        plot ({\xc + \rad*cos(\x)}, {\yc + \rad*sin(\x)});

    \filldraw[wordblue] ({\xc + \rad*cos(35)}, {\yc + \rad*sin(35)}) circle (2.5pt) node[below right, font=\footnotesize] {$R_\infty$ ($f \to \infty$)};
    \filldraw[wordblue] ({\xc + \rad*cos(145)}, {\yc + \rad*sin(145)}) circle (2.5pt) node[below left, font=\footnotesize] {$R_0 = I_0$ ($f \to 0$)};
    
    \filldraw[clinicred] ({\xc}, {\yc + \rad}) circle (3pt) node[above, font=\footnotesize\bfseries] {$-Z''_{\max} = \text{Max IP}$ ($f_c = \frac{1}{2\pi\tau}$)};
    \draw[dashed, gray] ({\xc}, 0) -- ({\xc}, {\yc + \rad});
    \draw[dashed, gray] (0, {\yc + \rad}) -- ({\xc}, {\yc + \rad});

    \draw[<->, thick, worddark] ({\xc + \rad*cos(145)}, -0.4) -- ({\xc + \rad*cos(35)}, -0.4) 
        node[midway, below, font=\footnotesize\bfseries] {$\Delta R = \text{DR} = R_0 - R_\infty$};

    \draw[dashed, gray] (\xc, \yc) -- (\xc+2, \yc);
    \draw[dashed, gray] (\xc, \yc) -- ({\xc + \rad*cos(35)}, {\yc + \rad*sin(35)});
    \draw[wordsteelblue, thick] (\xc+1, \yc) arc (0:35:1);
    \node[wordsteelblue, font=\footnotesize] at (\xc+1.5, \yc+0.3) {$\alpha\frac{\pi}{2}$};

    \fill[wordsteelblue!12, domain=35:145, samples=60] 
        plot ({\xc + \rad*cos(\x)}, {\yc + \rad*sin(\x)}) -- ({\xc + \rad*cos(35)},0) -- ({\xc + \rad*cos(145)},0) -- cycle;
    \node[wordblue, font=\small\bfseries] at (\xc, 1.2) {Área Espectral ($\text{Area} \propto \Delta R^2$)};
\end{tikzpicture}
}
\caption{Diagrama de Nyquist de Cole-Cole: Representación de los parámetros bioeléctricos extraídos por la plataforma a partir de la señal adquirida por la sonda VITA-PRB-0428.}
\label{fig:cole_cole}
\end{figure}

% ======================================================================================
% SECCIÓN 4: DATASETS CLÍNICOS Y FUSIÓN MULTIMODAL
% ======================================================================================
\section{Datasets Clínicos Internacionales y Fusión Multimodal Bayesiana}

\subsection{Los Tres Datasets de Referencia con Enlaces de Descarga}
En estricto apego a las directrices de ciencia de datos médicos, \textbf{los tres datasets no son combinados en una sola tabla plana}, pues corresponden a cohortes hospitalarias diferentes de Francia, Portugal y EE. UU. (Cuadro~\ref{tab:datasets}).

% --- CUADRO 2: DATASETS CLÍNICOS CON ENLACES (AJUSTADA A PÁGINA) ---
\begin{table}[H]
\centering
\small
\caption{Características de los tres datasets clínicos integrados en la plataforma VITA-ES.}
\label{tab:datasets}
\renewcommand{\arraystretch}{1.28}
\begin{tabularx}{\textwidth}{p{2.6cm} p{3.2cm} c p{4.3cm} c}
\rowcolor{tableheader}
\textcolor{white}{\textbf{Dataset}} & \textcolor{white}{\textbf{Cohorte y Repositorio}} & \textcolor{white}{\textbf{Muestras}} & \textcolor{white}{\textbf{Variables Clínicas Clave}} & \textcolor{white}{\textbf{Ponderación}} \\
\toprule
\textbf{1. EIS Mama} & J. Jossinet (1996) \cite{jossinet1996} \newline \href{https://archive.ics.uci.edu/dataset/192/breast+tissue}{\textcolor{wordsteelblue}{[UCI 192]}} \ $\bullet$ \ \href{https://www.kaggle.com/datasets/tarktunataalt/breast-tissue-impedance-measurements?resource=download}{\textcolor{wordsteelblue}{[Kaggle]}} & 106 ex-vivo & 9 descriptores Cole ($I_0$, $\text{PA500}$, $\text{DR}$, $\text{Max IP}$, $\text{Area}$) & \textbf{70\%} (Primaria) \\
\rowcolor{tablezebra}
\textbf{2. Coimbra BCCD} & M. Patrício et al. (2018) \cite{patricio2018} \newline \href{https://archive.ics.uci.edu/dataset/451/breast+cancer+coimbra}{\textcolor{wordsteelblue}{[UCI 451 BCCD]}} & 116 suero & Glucosa, Insulina, HOMA-IR, Leptina, Adiponectina, IMC & \textbf{15\%} (Metabólica) \\
\textbf{3. Wisconsin WDBC} & Wolberg et al. (1995) \cite{wolberg1995} \newline \href{https://archive.ics.uci.edu/dataset/17/breast+cancer+wisconsin+diagnostic}{\textcolor{wordsteelblue}{[UCI 17 WDBC]}} & 569 FNA & 30 variables nucleares morfométricas (Radio, Textura, Concavidad) & \textbf{15\%} (Morfológica) \\
\bottomrule
\end{tabularx}
\end{table}

\subsection{Parámetros Empíricos Centrales de la Cohorte Jossinet}
A partir del análisis directo de los 106 registros de biopsia confirmada de \texttt{breast\_tissue.csv}, se calibraron las distribuciones empíricas ($\mu \pm \sigma$) utilizadas en el motor de inferencia (Cuadro~\ref{tab:empirical_stats}).

% --- CUADRO 3: PARÁMETROS EMPÍRICOS (AJUSTADA CON RESIZEBOX) ---
\begin{table}[H]
\centering
\small
\caption{Parámetros empíricos biofísicos extraídos del dataset de bioimpedancia mamaria ($N=106$).}
\label{tab:empirical_stats}
\renewcommand{\arraystretch}{1.25}
\resizebox{\columnwidth}{!}{
\begin{tabular}{l c c c c c}
\rowcolor{tableheader}
\textcolor{white}{\textbf{Tejido Mamario (Clase)}} & \textcolor{white}{\textbf{$I_0$ ($\Omega$)}} & \textcolor{white}{\textbf{$\text{DR}$ ($\Omega$)}} & \textcolor{white}{\textbf{Max IP ($\Omega$)}} & \textcolor{white}{\textbf{$\text{Area}$ ($\Omega^2$)}} & \textcolor{white}{\textbf{Alpha ($\alpha$)}} \\
\toprule
\textbf{Carcinoma (\texttt{car}, $n=21$)} & $394.2 \pm 87.1$ & $153.0 \pm 67.0$ & $64.5 \pm 18.6$ & $5,723.1 \pm 3,172.3$ & $0.85 \pm 0.18$ \\
\rowcolor{tablezebra}
\textbf{Fibroadenoma (\texttt{fad}, $n=15$)} & $245.9 \pm 69.8$ & $47.8 \pm 22.3$ & $23.0 \pm 10.9$ & $610.2 \pm 430.6$ & $0.82 \pm 0.16$ \\
\textbf{Mastopatía (\texttt{mas}, $n=18$)} & $290.3 \pm 112.0$ & $63.6 \pm 39.2$ & $31.5 \pm 10.8$ & $1,082.4 \pm 1,254.1$ & $0.85 \pm 0.17$ \\
\rowcolor{tablezebra}
\textbf{Glandular (\texttt{gla}, $n=16$)} & $238.3 \pm 119.2$ & $27.5 \pm 33.2$ & $26.4 \pm 7.2$ & $411.7 \pm 632.8$ & $0.96 \pm 0.06$ \\
\textbf{Conectivo (\texttt{con}, $n=14$)} & $1,212.9 \pm 386.5$ & $358.0 \pm 179.6$ & $73.0 \pm 34.5$ & $5,324.6 \pm 4,069.2$ & $0.57 \pm 0.14$ \\
\rowcolor{tablezebra}
\textbf{Adiposo (\texttt{adi}, $n=22$)} & $2,052.1 \pm 342.5$ & $324.7 \pm 230.4$ & $194.6 \pm 106.6$ & $24,889.7 \pm 35,695.3$ & $0.88 \pm 0.16$ \\
\bottomrule
\end{tabular}
}
\end{table}

\subsection{Ecuación de Fusión Multimodal Bayesiana}
La probabilidad combinada de malignidad ($P_{\text{Fused}}$) sintetiza los tres canales de evidencia:

\begin{equation}
    P_{\text{Fused}} = 0.70 \cdot P_{\text{car}}(\text{EIS}) + 0.15 \cdot P_{\text{Metabólico}}(\text{Coimbra}) + 0.15 \cdot P_{\text{Morfológico}}(\text{WDBC})
    \label{eq:fused}
\end{equation}

% --- FIGURA 3: FLUJO DE FUSIÓN MULTIMODAL EN TIKZ (RESIZED) ---
\begin{figure}[H]
\centering
\resizebox{0.90\linewidth}{!}{
\begin{tikzpicture}[
    node distance=1.1cm,
    box/.style={rectangle, draw=wordblue, fill=boxbg, rounded corners=3pt, text centered, font=\small, minimum height=1.0cm},
    box_model/.style={rectangle, draw=worddark, fill=white, rounded corners=3pt, text centered, font=\footnotesize, minimum height=0.9cm},
    fusion/.style={ellipse, draw=huaweired, fill=huaweired!10, text centered, font=\small\bfseries, minimum height=1.1cm},
    triage/.style={rectangle, draw=clinicred, fill=clinicred!8, rounded corners=4pt, text centered, font=\small\bfseries, minimum height=1.2cm},
    arrow/.style={-{Latex[length=2.5mm]}, thick, worddark}
]
    \node (in_eis) [box, text width=3.8cm] {\textbf{Sonda VITA-PRB-0428} \\ Barrido EIS (15 kHz--1 MHz)};
    \node (in_coim) [box, right=0.7cm of in_eis, text width=3.8cm] {\textbf{Perfil Metabólico} \\ Glucosa, Insulina, HOMA, IMC};
    \node (in_wdbc) [box, right=0.7cm of in_coim, text width=3.8cm] {\textbf{Hallazgo Palpatorio} \\ Presencia de Nódulo / Atipia};

    \node (m_eis) [box_model, below=0.8cm of in_eis, text width=3.8cm] {\textbf{Model Zoo EIS} \\ XGBoost / VITA-Net 1D \\ $P_{\text{car}}$ ($w_1 = 0.70$)};
    \node (m_coim) [box_model, below=0.8cm of in_coim, text width=3.8cm] {\textbf{Modelo Coimbra} \\ Random Forest Calibrado \\ $P_{\text{Met}}$ ($w_2 = 0.15$)};
    \node (m_wdbc) [box_model, below=0.8cm of in_wdbc, text width=3.8cm] {\textbf{Modelo Morfológico} \\ Regresión Logística WDBC \\ $P_{\text{Morf}}$ ($w_3 = 0.15$)};

    \node (fuse) [fusion, below=0.9cm of m_coim, text width=6.5cm] {\textbf{Fusión Bayesiana Tardía} \\ $P_{\text{Fused}} = 0.70 P_{\text{car}} + 0.15 P_{\text{Met}} + 0.15 P_{\text{Morf}}$};

    \node (triage_res) [triage, below=0.8cm of fuse, text width=11.0cm] {
        \textbf{Dictamen Clínico de Triaje VITA-ES} \\
        \footnotesize $\bullet$ $P_{\text{Fused}} \ge 32\%$ $\to$ \textbf{ALTO RIESGO (BI-RADS 4/5)} | Derivación Biopsia Urgente ($<7$ días) \\
        \footnotesize $\bullet$ $P(\text{fad})+P(\text{mas}) \ge 38\%$ $\to$ \textbf{RIESGO MODERADO (BI-RADS 3)} | Ecografía a 3--6 meses \\
        \footnotesize $\bullet$ Fisiológico $\to$ \textbf{BAJO RIESGO (BI-RADS 1/2)} | Control Anual Preventivo
    };

    \draw [arrow] (in_eis) -- (m_eis);
    \draw [arrow] (in_coim) -- (m_coim);
    \draw [arrow] (in_wdbc) -- (m_wdbc);
    \draw [arrow] (m_eis) |- (fuse);
    \draw [arrow] (m_coim) -- (fuse);
    \draw [arrow] (m_wdbc) |- (fuse);
    \draw [arrow] (fuse) -- (triage_res);

\end{tikzpicture}
}
\caption{Flujo de Fusión Multimodal Tardía: Integración probabilística de impedancia celular, perfil metabólico en sangre y morfología física.}
\label{fig:fusion}
\end{figure}

% ======================================================================================
% SECCIÓN 5: MODEL ZOO E INCERTIDUMBRE
% ======================================================================================
\section{Model Zoo Clínico y Cuantificación de Incertidumbre (CUI)}

\subsection{Benchmarking de los Clasificadores (Stratified $k=5$)}
Se implementó un Model Zoo comparativo para equilibrar sensibilidad y latencia (Cuadro~\ref{tab:model_zoo}). En medicina preventiva de tamizaje, la métrica crítica es la \textbf{Sensibilidad en Carcinoma} (minimización estricta de falsos negativos).

% --- CUADRO 4: MODEL ZOO (AJUSTADO CON RESIZEBOX PARA EVITAR DESBORDAMIENTO) ---
\begin{table}[H]
\centering
\small
\caption{Rendimiento comparativo del Model Zoo clínico VITA-ES ($k=5$ Stratified Cross-Validation).}
\label{tab:model_zoo}
\renewcommand{\arraystretch}{1.25}
\resizebox{\columnwidth}{!}{
\begin{tabular}{l c c c c c}
\rowcolor{tableheader}
\textcolor{white}{\textbf{Modelo Clasificador}} & \textcolor{white}{\textbf{Exactitud (\%)}} & \textcolor{white}{\textbf{Bal. Acc. (\%)}} & \textcolor{white}{\textbf{F1 Macro}} & \textcolor{white}{\textbf{Sensibilidad Ca. (\%)}} & \textcolor{white}{\textbf{Latencia Edge (ms)}} \\
\toprule
\textbf{Balanced Random Forest} & $77.36 \pm 4.2$ & $78.43$ & $0.7612$ & $85.71$ & $4.18$ \\
\rowcolor{tablezebra}
\textbf{SVM (Kernel RBF)} & $73.58 \pm 5.1$ & $74.20$ & $0.7250$ & $80.95$ & $\mathbf{0.17}$ \\
\textbf{XGBoost (Campeón)} & $\mathbf{83.02} \pm 3.8$ & $\mathbf{84.55}$ & $\mathbf{0.8415}$ & $\mathbf{95.24}$ & $2.66$ \\
\rowcolor{tablezebra}
\textbf{VITA-Net 1D (CNN)} & $78.30 \pm 4.9$ & $79.12$ & $0.7780$ & $90.48$ & $5.12$ \\
\bottomrule
\end{tabular}
}
\end{table}

\subsection{Formulación Matemática del Clinical Uncertainty Index (CUI)}
Para prevenir fallos catastróficos por sobreconfianza ante atipias biológicas limítrofes, el sistema evalúa el índice CUI a partir del vector de distribución de probabilidades $\mathbf{p} \in \mathbb{R}^6$:

\begin{equation}
    H_{\text{norm}} = \frac{-\sum_{i=1}^6 p_i \ln p_i}{\ln 6}, \quad \Delta P = P_{(1)} - P_{(2)}
\end{equation}

\begin{equation}
    \text{CUI} = \Big[ 0.65 \cdot H_{\text{norm}} + 0.35 \cdot (1 - \Delta P) \Big] \times 100\%
    \label{eq:cui}
\end{equation}

\begin{itemize}
    \item \textbf{CUI $\le 35\%$ (Certeza Alta):} Distribución unimodal pronunciada. Se autoriza el dictamen de triaje automático.
    \item \textbf{CUI $36\% - 55\%$ (Certeza Moderada):} Ambigüedad limítrofe entre clases benignas (p. ej., fibroadenoma vs. mastopatía). Se recomienda ecografía complementaria.
    \item \textbf{CUI $> 55\%$ (Incertidumbre Crítica):} Interlock de seguridad activo. Se suspende el descarte preventivo y se emite la recomendación obligatoria: \textit{«Muestra no concluyente: Requiere confirmación ecográfica/histológica»}.
\end{itemize}

% ======================================================================================
% SECCIÓN 6: MARCO REGULATORIO EN EL SALVADOR
% ======================================================================================
\section{Marco Regulatorio y Ruta de Validación Clínica (SRS El Salvador)}

Dada la experiencia regulatoria del equipo en la **Superintendencia de Regulación Sanitaria (SRS)** de El Salvador, el desarrollo de VITA-ES se ciñe estrictamente a la normativa legal vigente para dispositivos médicos:

\subsection{Normativa Local Salvadoreña}
\begin{enumerate}
    \item \textbf{Clasificación del Dispositivo (Clase II):} Conforme al **RTS 11.03.02:21** (\textit{Reglamento Técnico Salvadoreño para el Registro de Dispositivos Médicos}), VITA-ES clasifica como Dispositivo Médico Clase II (riesgo potencial moderado, no invasivo con fuente de energía externa para tamizaje diagnóstico).
    \item \textbf{Buenas Prácticas de Manufactura y Ensamble:} Cumplimiento del **RTS 11.03.03:22** para ensamble de sondas, almacenamiento y distribución de parches biocompatibles.
    \item \textbf{Comité de Ética en Investigación:} Aprobación previa por el **Comité Nacional de Ética de Investigación en Salud (CNEIS)** y comités locales (CLEIS) para la ejecución de ensayos con pacientes.
    \item \textbf{Tecnovigilancia y Registro Sanitario:} Trazabilidad y reporte continuo de eventos adversos bajo el **RTS 11.03.04:23**.
\end{enumerate}

\subsection{Estándares Internacionales de Ingeniería Biomédica}
\begin{itemize}
    \item \textbf{ISO 13485:2016:} Sistema de gestión de la calidad para diseño y manufactura de dispositivos médicos.
    \item \textbf{ISO 14971:2019:} Gestión de riesgos para dispositivos médicos a lo largo de todo el ciclo de vida.
    \item \textbf{IEC 60601-1 / IEC 60601-1-2:} Seguridad eléctrica y compatibilidad electromagnética en equipos electromédicos.
    \item \textbf{IEC 62304 / IEC 62366:} Ciclo de vida del software médico (SaMD) e ingeniería de usabilidad para enfermería comunitaria.
\end{itemize}

\subsection{Fases de Validación Clínica (Protocolo)}
\begin{itemize}
    \item \textbf{Fase 0 (Preclínica de Laboratorio):} Validación biofísica con fantomas de gel conductivo mamario (coeficiente de variación $CV < 5\%$).
    \item \textbf{Fase 1 (Piloto Clínico en UCSF):} Estudio con $N=100$ pacientes en unidades de salud comunitarias contrastado contra mamografía digital bilateral.
    \item \textbf{Fase 2 (Validación Multicéntrica):} $N=500$ pacientes en 5 departamentos de El Salvador para re-entrenamiento continuo de los modelos de IA en Huawei ModelArts con fenotipos mamarios locales.
    \item \textbf{Fase 3 (Alcance y Financiamiento):} Registro sanitario formal ante la SRS y postulación a financiamiento mediante subsidio **R21 de los NIH** (\textit{National Institutes of Health} de EE. UU.) a través de colaboración institucional internacional.
\end{itemize}

% ======================================================================================
% SECCIÓN 7: ARQUITECTURA HUAWEI CLOUD
% ======================================================================================
\section{Arquitectura Tecnológica Huawei Cloud (Edge-to-Cloud)}

VITA-ES implementa una topología distribuida que articula la inferencia ultrarrápida en la sonda portátil con la gobernanza hospitalaria y re-entrenamiento continuo en la nube de Huawei.

% --- FIGURA 4: ARQUITECTURA HUAWEI CLOUD EN TIKZ (COMPACTA Y SIN DESBORDES) ---
\begin{figure}[H]
\centering
\resizebox{0.95\linewidth}{!}{
\begin{tikzpicture}[
    node distance=0.8cm,
    hw_edge/.style={rectangle, draw=wordblue, fill=wordblue!6, rounded corners=4pt, text centered, font=\small, minimum height=0.9cm},
    hw_cloud/.style={rectangle, draw=huaweired, fill=huaweired!6, rounded corners=4pt, text centered, font=\small, minimum height=0.9cm},
    subnode/.style={rectangle, draw=worddark!40, fill=white, rounded corners=2pt, text centered, font=\footnotesize, minimum height=0.7cm},
    arrow_t/.style={-{Latex[length=2.5mm]}, thick, worddark}
]
    % DISPOSITIVO EDGE
    \node (sonda) [hw_edge, text width=6.2cm] {\textbf{SONDA PORTÁTIL VITA-PRB-0428} \\ AD5933 (12-bit) + ESP32-S3 BLE 5.0 (36.8~$^\circ$C, LiPo 94\%)};
    \node (edge_ai) [subnode, below=0.6cm of sonda, text width=5.8cm] {\textbf{Huawei MindSpore Lite (Ascend 310)} \\ Inferencia Local en Sonda ($4.2\text{ ms}$, Offline Ready)};

    % ENLACE TELEMETRÍA
    \draw [arrow_t] (sonda) -- (edge_ai);

    % CLOUD GATEWAY (ESPACIAMIENTO CONTROLADO)
    \node (ecs) [hw_cloud, right=1.0cm of sonda, text width=6.2cm] {\textbf{Huawei ECS Gateway (EulerOS c7.large.2)} \\ API Microservicios RESTful / Validación SQI y TLS 1.3};

    % SERVICIOS CLOUD
    \node (modelarts) [subnode, below=0.5cm of ecs, text width=5.8cm] {\textbf{Huawei ModelArts (Ascend 910 Pro)} \\ Inferencia Online Serverless / MLOps Federado};
    \node (dcs) [subnode, below=0.45cm of modelarts, text width=5.8cm] {\textbf{Huawei DCS (Redis Enterprise)} \\ Buffer de Telemetría Streaming (Latencia $< 1.5\text{ ms}$)};
    \node (rds) [subnode, below=0.45cm of dcs, text width=5.8cm] {\textbf{Huawei RDS (PostgreSQL 14 Cifrado)} \\ Expedientes KMS AES-256 (HIPAA/GDPR Compliant)};

    % Conexiones Edge-Cloud
    \draw [arrow_t, dashed] (sonda.east) -- node[above, font=\footnotesize\bfseries] {GATT / 4G} (ecs.west);
    \draw [arrow_t] (ecs) -- (modelarts);
    \draw [arrow_t] (modelarts) -- (dcs);
    \draw [arrow_t] (dcs) -- (rds);

    % Cajas de agrupación
    \begin{scope}[on background layer]
        \node [draw=wordblue, fill=wordblue!2, fit=(sonda)(edge_ai), inner sep=0.25cm, rounded corners=5pt, label={[wordblue, font=\bfseries]above:DISPOSITIVO EDGE (LOCAL / SIN RED)}] {};
        \node [draw=huaweired, fill=huaweired!2, fit=(ecs)(modelarts)(dcs)(rds), inner sep=0.25cm, rounded corners=5pt, label={[huaweired, font=\bfseries]above:HUAWEI CLOUD SUITE}] {};
    \end{scope}

\end{tikzpicture}
}
\caption{Arquitectura Edge-to-Cloud de VITA-ES: Sincronización entre la sonda portátil y los servicios de Huawei Cloud.}
\label{fig:huawei_arch}
\end{figure}

\subsection{Servicios de Huawei Cloud Desplegados}
\begin{enumerate}
    \item \textbf{Huawei ModelArts (Inferencia y MLOps):} Empaquetado en \texttt{modelarts\_deploy/} con scripts de servicio \texttt{customize\_service.py} y descriptores \texttt{config.json}. Permite orquestación de modelos en contenedores acelerados por NPU Huawei Ascend 910 Pro con re-entrenamiento federado continuo para la Fase 2 del estudio multicéntrico.
    \item \textbf{MindSpore Lite Edge Ingestion:} Exportación de modelos a formato abierto ONNX (\texttt{model\_rf.onnx} y \texttt{model\_svm.onnx}) para ejecución en la sonda autónoma con procesadores Ascend 310, garantizando operatividad en brigadas rurales sin internet.
    \item \textbf{Huawei DCS (Distributed Cache Service):} Clúster Redis Enterprise para almacenar tramas binarias BLE crudas (\texttt{AA 55 04 28 ...}) a latencias inferiores a 1.5~ms, amortiguando ráfagas en jornadas de tamizaje masivo.
    \item \textbf{Huawei RDS (Relational Database Service):} Base de datos relacional PostgreSQL con cifrado en reposo administrado mediante \textbf{Huawei KMS (Key Management Service)} con estándar AES-256 para salvaguardar expedientes clínicos conforme a normativas HIPAA, GDPR y RTS 11.03.04:23.
\end{enumerate}

% ======================================================================================
% REFERENCIAS BIBLIOGRÁFICAS
% ======================================================================================
\begin{thebibliography}{99}
\bibitem{jossinet1996}
J. Jossinet, ``Variability of impedivity in normal and pathological breast tissue,'' \textit{Medical \& Biological Engineering \& Computing}, vol. 34, no. 5, pp. 346--350, 1996.

\bibitem{uci_breast_tissue}
J. Jossinet, ``Breast Tissue Electrical Impedance Measurements Dataset,'' \textit{UCI Machine Learning Repository}, 2010. Disponible en: \url{https://archive.ics.uci.edu/dataset/192/breast+tissue}

\bibitem{kaggle_breast_tissue}
T. Tunataalt, ``Breast Tissue Impedance Measurements (Ex-Vivo Multiclass Spectral Data),'' \textit{Kaggle Datasets}, 2021. Disponible en: \url{https://www.kaggle.com/datasets/tarktunataalt/breast-tissue-impedance-measurements?resource=download}

\bibitem{patricio2018}
M. Patrício, J. Pereira, J. Crisóstomo, P. Matafome, M. Gomes, R. Seiça, and F. Caramelo, ``Using Resistin, Glucose, Age and BMI to Predict the Presence of Breast Cancer,'' \textit{BMC Cancer}, vol. 18, no. 1, p. 29, 2018.

\bibitem{uci_coimbra}
M. Patrício et al., ``Breast Cancer Coimbra (BCCD) Clinical Serum Dataset,'' \textit{UCI Machine Learning Repository}, 2018. Disponible en: \url{https://archive.ics.uci.edu/dataset/451/breast+cancer+coimbra}

\bibitem{wolberg1995}
W. H. Wolberg, W. N. Street, and O. L. Mangasarian, ``Breast Cancer Wisconsin (Diagnostic) Data Set,'' \textit{UCI Machine Learning Repository}, 1995. Disponible en: \url{https://archive.ics.uci.edu/dataset/17/breast+cancer+wisconsin+diagnostic}

\bibitem{cole1941}
K. S. Cole and R. H. Cole, ``Dispersion and absorption in dielectrics I. Alternating current characteristics,'' \textit{The Journal of Chemical Physics}, vol. 9, no. 4, pp. 341--351, 1941.

\bibitem{silva2000}
J. E. Silva et al., ``Sensitivity and specificity of electrical impedance spectroscopy in breast cancer screening,'' \textit{IEEE Transactions on Biomedical Engineering}, vol. 47, no. 6, pp. 820--827, 2000.

\bibitem{acr2013}
American College of Radiology (ACR), \textit{ACR BI-RADS Atlas: Breast Imaging Reporting and Data System}, 5th ed. Reston, VA: American College of Radiology, 2013.

\bibitem{analog2020}
Analog Devices, \textit{AD5933: 1 MSPS, 12-Bit Impedance Converter Network Analyzer Technical Datasheet}, Rev. F, 2020.

\bibitem{huawei2025}
Huawei Cloud, \textit{ModelArts AI Development Platform \& MindSpore Lite Edge Inference Architecture Guide}, Huawei Technologies Co., Ltd., 2025.

\bibitem{srs2021}
Superintendencia de Regulación Sanitaria (SRS) de El Salvador, \textit{RTS 11.03.02:21 Dispositivos Médicos. Registro Sanitario e Importación}, San Salvador, El Salvador, 2021.

\bibitem{srs2023}
Superintendencia de Regulación Sanitaria (SRS) de El Salvador, \textit{RTS 11.03.04:23 Ensayos Clínicos con Dispositivos Médicos y Tecnovigilancia}, San Salvador, El Salvador, 2023.
\end{thebibliography}

\end{document}
